\hspace*{\fill} %标题前空两行

\section{绪论}

\subsection{研究背景与意义}

路径规划是计算机科学里的一个基础问题，应用范围很广——机器人导航、游戏NPC行为设计、物流配送路线优化、网络数据包转发，背后都是找路的问题。在游戏开发领域，路径规划算法的质量直接决定游戏AI的表现。一个总是在墙角卡死的怪物会瞬间毁掉玩家的沉浸感，而一个能流畅绕行的敌人则能给玩家带来真正的挑战。

贪吃蛇游戏规则很简单——控制蛇吃食物、变长、避免撞墙和撞自身——但它是一个很有意思的算法验证场景。游戏的状态空间是动态变化的，蛇每吃一个食物，身体就增长一格，可用空间就缩小一些。这意味着同一个路径规划算法在游戏开局和后期的表现可能截然不同。比起静态迷宫问题，贪吃蛇更考验算法的自适应能力。

选A*算法作为核心寻路方案，因为它在理论完备性和实际效率之间找到了不错的平衡点\upcite{2}。A*最早由Hart等人在1968年提出\upcite{1}，通过引入启发式函数引导搜索方向，避免了盲目穷举的巨大计算开销。但这个效率不是白来的——启发式函数的设计需要对具体问题有一定了解。选得好，搜索又快又准；选得不好，要么搜索不充分漏掉好路径，要么退化成暴力搜索。这种代价在实际开发中是躲不开的。

用Python和Pygame实现这个系统，主要是考虑开发效率。Python语法简洁，加上Pygame的2D图形渲染能力，可以快速搭出可视化原型。Python比C/C++慢很多，这是事实，但对贪吃蛇这种网格规模不大的游戏来说，性能瓶颈不算严重。如果以后要扩展到更大规模的网格或更复杂的场景，可以考虑用Cython或迁移到C++。

\subsection{国内外研究现状}

A*算法自提出以来经历了数十年发展，衍生出大量变体和优化方案。理论上，原始A*是完备且最优的——只要启发式函数是可采纳的（admissible），即从不高估到达目标的实际代价，A*就能保证找到最短路径\upcite{1}。但"最优"只是理论保证，实际应用中受内存和时间限制，完全展开所有节点往往不可行。

在网格路径规划领域，Sturtevant对多种搜索算法做了系统的基准测试\upcite{6}，发现在规则网格上A*的性能与启发式函数的选择高度相关。Manhattan距离在四方向移动的网格中表现很好，但换成八方向或允许对角线移动时，它的可采纳性就被打破了，需要用切比雪夫距离或Octile距离替代。这个细节在很多教程里轻描淡写，但实际踩坑的时候才会意识到它的重要性。

在游戏工业界，A*算法的应用更广泛，但也面临更多工程层面的挑战\upcite{5}。Cui和Shi在2011年的研究中指出，现代计算机游戏中A*寻路的瓶颈往往不在算法本身，而在大规模图的存储和节点扩展效率上\upcite{5}。针对这个问题，游戏开发者引入了跳点搜索（Jump Point Search, JPS）、分层路径规划（Hierarchical Pathfinding）等优化技术。JPS通过识别对称路径中的"跳点"大幅减少需要扩展的节点数量，在均匀代价的网格上效果显著。但它对启发式函数要求更高，在非均匀代价网格上的适用性有限。

国内学者在路径规划算法方面也做了不少工作。王晓东在《计算机算法设计与分析》中系统介绍了A*算法及其变体的理论基础\upcite{7}，并对不同搜索策略的时间复杂度和空间复杂度做了详尽比较。张铭等人在教材中侧重于算法工程实现，给出了较为清晰的伪代码框架\upcite{10}。不过教材里的例子大多偏向理论验证，和真实游戏场景中的动态环境、实时响应需求之间还是有距离的。

在贪吃蛇AI这个具体方向上，社区方案五花八门。有人用简单的贪心策略——每步都朝食物走最近的路——效果在游戏前期不错，但蛇一长就容易把自己逼死。也有人用Hamilton回路来保证蛇永远不死，但这需要预设一条遍历所有格子的固定路径，蛇的行为看起来毫无智能可言，而且构造Hamilton回路本身也需要额外计算。A*算法则介于两者之间：不像纯贪心那样短视，也不像Hamilton回路那样僵化。代价也有——在蛇身很长、空间逼仄的时候，A*可能找不到安全路径，需要配合"逃生模式"等策略兜底。

LaValle在规划算法领域的工作\upcite{9}提供了更广视角的理论框架，虽然他的研究更多面向连续空间中的运动规划，但其中关于搜索完备性、最优性、计算复杂度的讨论对离散网格上的路径规划同样有启发意义。

\subsection{主要研究内容}

本文的主要工作围绕以下几个方面展开：

第一，A*算法的原理分析与实现。从Hart等人提出的原始公式出发\upcite{1}，理解$f(n)=g(n)+h(n)$中代价函数$g(n)$和启发式函数$h(n)$各自的作用，以及它们之间的权重关系对搜索行为的影响。实现时需要处理好开放列表（open list）和关闭列表（closed list）的数据结构选择——Python自带的列表在频繁插入删除操作下性能堪忧，用heapq模块实现的优先队列是更合理的选择\upcite{4}。

第二，启发式函数的设计与对比。针对贪吃蛇游戏的四方向移动网格，选择Manhattan距离作为基础启发式函数，并探讨不同权重系数对搜索性能的影响。这部分工作量不大，但反复调参的过程确实让人体会到"算法性能对参数敏感"这句话不是论文里的套话。

第三，游戏系统的整体设计与开发。基于Pygame框架搭建贪吃蛇游戏的基本框架，包括游戏主循环、事件处理、碰撞检测、渲染绘制等模块。在此基础上集成A*寻路模块，实现手动操控与AI自动运行两种模式。

第四，系统测试与性能分析。在不同网格规模和游戏阶段下对A*算法的路径规划效果进行测试，记录搜索节点数、路径长度、单帧耗时等指标，分析算法在蛇身增长导致空间变窄时的表现退化情况。

\subsection{论文组织结构}

本文共分为六章。

第一章介绍研究背景、国内外研究现状和主要研究内容。

第二章介绍相关技术基础，包括Python语言与Pygame框架的特点、A*寻路算法的原理、BFS与Dijkstra算法的对比分析、启发式函数的设计考量，以及游戏AI的决策机制。

第三章进行系统需求分析与总体设计，从功能需求和非功能需求两个维度出发，给出系统的架构设计、模块划分和数据流设计。

第四章详细描述系统各模块的具体实现过程，包括游戏核心逻辑、A*算法实现、AI决策策略和界面交互等。

第五章对系统进行测试与分析，给出不同场景下的测试结果，并对算法性能进行量化评估。

第六章总结全文工作，分析当前系统的不足之处，展望可能的改进方向。
